% ch04 —— §4 相互局域修正（原文 wcf-paper.tex 1389–2128 行，对应 PDF 第 22–34 页）
% 原文宏 \plog、\eps 不在本工程 preamble 中，此处本地补齐（定义与原文 wcf-paper.tex 159、166 行相同）。
% 原文 (4.42) 中一处 \biggr) 疑为 \right) 之笔误（原文 1720 行），转录保持原文，编译无碍。
\providecommand{\plog}[1]{#1\partial_{#1}}
\providecommand{\eps}{\epsilon}

\section{相互局域修正}
\label{sec:periodicnut}

如果我们只考虑无质量部分的朴素量纲约化，那么半平坦度量 $g^\sfm$ 就是全部答案了。然而，$d=4$ 的理论中还有有质量的 BPS 粒子。当这些有质量粒子中的一个或多个绕 $S^1$ 转一圈时，相应的「瞬子」组态会给度量带来修正。由于 $d=4$ 理论中状态的能量满足 $M \ge \abs{Z}$，这些修正的权重至少是 $e^{-2 \pi R \abs{Z}}$。

本节所研究的例子中，所有 BPS 粒子都相互局域。这比一般情形容易处理得多，因为我们可以选一个对偶标架，使这些粒子全都带电荷，于是完全可以只在有效拉格朗日描述之内工作。

本节大部分篇幅里，我们进一步特化到耦合单个带电荷超多重态的自由 $U(1)$ 规范理论。这个理论既是最简单的例子，也在物理上切题：它描述的是 $\vm$ 中一般奇点附近的物理——在那里，有一个 BPS 粒子变得比其余粒子轻得多。

\subsection{精确的单粒子度量}

我们考虑 $\IR^3 \times S^1$ 上耦合单个电荷为 $q>0$ 的超多重态（连同其电荷为 $-q$ 的 CPT 共轭）的 $U(1)$ 规范理论。我们将要描述的度量，此前已在 \cite{Ooguri:1996me,Seiberg:1996ns} 中考虑过。

$d=4$ 理论的模空间 $\vm$ 由矢量多重态标量 $a \in \IC$ 坐标化。更确切地说，$\vm$ 只是 $\IC$ 中的一个开片，因为 $d=4$ 理论不是渐近自由的：在 $\abs{a} \sim \abs{\Lambda}$ 处存在截断。

正如我们在第 2.3 节中所解释的，$d=3$ 理论的模空间 $\CM$ 是 $\vm$ 上的 2-环面丛。环面纤维 $\CM_a$ 由电 Wilson 线 $\theta_\elec$ 与磁 Wilson 线 $\theta_\magn$ 坐标化（这里暂时忽略二次加细（quadratic refinement）的微妙之处），两者的周期都是 $2 \pi$。% TODO 合稿改 \ref{sec:reduction}

半平坦度量 $g^\sfm$ 具有 $U(1)^2$ 的等距作用，因为 $\theta_\elec$ 与 $\theta_\magn$ 的平移都是精确对称性。带电荷的超多重态与 $\theta_\elec$ 耦合，因而破坏了平移 $\theta_\elec$ 的那个等距。不过理论中没有带磁荷的 BPS 态，所以 $\theta_\magn$ 的平移仍是精确等距。因此修正后的度量 $g$ 具有 Gibbons–Hawking 形式。

为与 Gibbons–Hawking 拟设作比较，我们引入矢量 $\vec{x}$：
\begin{equation}
a = x^1 + i x^2, \quad \theta_\elec = 2 \pi R x^3.
\end{equation}
$\theta_\magn$ 是由 $\vec{x}$ 参数化的 $\IR^2 \times S^1$ 开子集上某个 $U(1)$ 丛的局部坐标。度量为
\begin{equation}
g = V(\vec{x})^{-1} \left(\frac{d \theta_\magn}{2\pi} + A(\vec{x})\right)^2 + V(\vec{x}) d\vec{x}^2,
\end{equation}
其中 $V$ 是正的调和函数（下文将算出它），$A$ 是曲率为
\begin{equation} \label{eq:FdV}
F = \star dV.
\end{equation}
的 $U(1)$ 联络。这是标准 Gibbons–Hawking 拟设的一个轻微推广：标准拟设中取 $\vec{x}$ 落在 $\IR^3$（的某个开子集）里。（我们可以先在 $\IR^3$ 的合适子集上工作，再对总空间上平移 $x^3$ 的 $\IZ$-作用取商。）在标准拟设中，所有满足 \eqref{eq:FdV} 的 $A$ 都互相规范等价，因而定义同一个度量。在我们的情形中这并不完全成立：与绕 $S^1$ 的 holonomy 相关，还有一个额外的规范不变自由度。这一选择与 $d=4$ 中 $\theta$ 角的选择有关。

在我们的情形中，$V(\vec{x})$ 可以通过一圈计算积分掉带电荷超多重态来得到。文献 \cite{Seiberg:1996ns} 断言了一个非重整化定理，它意味着这一计算是精确的。所得的 $V$ 是 $\IR^2 \times S^1$ 上带 $q$ 个奇点的调和函数。$\theta_\elec$ 的周期性来自对带电荷超多重态在 $S^1$ 上的 Kaluza–Klein 动量求和：
\begin{equation} \label{eq:V-sum}
V = \frac{q^2 R}{4 \pi} \sum_{n=-\infty}^\infty \left( \frac{1}{\sqrt{q^2 R^2 \abs{a}^2 + (q \frac{\theta_\elec}{2 \pi} + n)^2}} - \kappa_n \right)
\end{equation}
这里 $\kappa_n$ 是为使求和收敛而引入的正规化常数。\eqref{eq:V-sum} 的 Poisson 重求和表明
\begin{equation}
V = V^\sfm + V^\inst,
\end{equation}
其中
\begin{align}
V^\sfm &= -\frac{q^2 R}{4\pi} \left(\log  \frac{a}{\Lambda} + \log \frac{\bar{a}}{\bar{\Lambda}} \right), \\
V^\inst &= \frac{q^2 R}{2\pi} \sum_{n\neq 0} e^{i n q \theta_\elec} K_0(2 \pi R \abs{n q a}). \label{eq:V-inst}
\end{align}
这里 $\Lambda$ 是与 $\kappa_n$ 的选取有关的紫外截断。\footnote{例如，若取 $\kappa_n = (\vert \tilde \Lambda\vert^2 + n^2)^{-1/2}$，则可取 $\Lambda = (qR)^{-1} \tilde\Lambda \exp[-2 \sum_{m=1}^\infty K_0(2\pi m \vert \tilde \Lambda\vert )]$。}

要完全指定度量，还须给出满足 \eqref{eq:FdV} 的 $A(\vec{x})$：
\begin{equation}
A = A^\sfm + A^\inst,
\end{equation}
其中
\begin{align}
A^\sfm &= \frac{i q^2}{8 \pi^2} \left( \log \frac{a}{\Lambda} - \log \frac{\bar{a}}{\bar{\Lambda}} \right) d\theta_\elec, \label{eq:A0} \\
A^\inst &=  - \frac{q^2 R}{4 \pi} \left( \frac{da}{a} - \frac{d \bar a}{\bar a} \right) \sum_{n \neq 0} (\sgn n) e^{i n q \theta_\elec} \abs{a} K_1(2 \pi R \abs{n q a}).  \label{eq:A-inst}
\end{align}
在大 $R$ 时，$V$ 与 $A$ 中的领头项是 $V^\sfm$ 与 $A^\sfm$。只保留这些项，$g$ 就成为半平坦度量，其
\begin{equation}
\tau = \frac{q^2}{2 \pi i} \log \frac{a}{\Lambda}.
\end{equation}
这就是 $d=4$ 中积分掉超多重态所得的跑动耦合。

次领头项 $V^\inst$、$A^\inst$ 给出对半平坦度量的修正。由于 $x \to +\infty$ 时的渐近行为 $K_\nu(x) \sim \sqrt{\frac{\pi}{2x}} e^{-x}$，它们如我们所预期地呈瞬子展开的形式。它们也如预期地破坏了沿 $\theta_\elec$ 的平移不变性。最后，它们改善了奇性。回忆在 $g^\sfm$ 中 $a=0$ 处有一个奇点。由 \eqref{eq:V-sum} 可见，$g$ 唯一可能的奇点出现在 $a = 0$、$q \theta_\elec = 2 \pi n$ 处。研究这些点附近的度量可知，每一处都是 $A_{q-1}$ 型锥奇点。于是 $g^\sfm$ 中的一个奇点被 $g$ 中的 $q$ 个更高余维的奇点取代。在最简单的 $q=1$ 情形，奇性被完全光滑化。

\subsubsection*{整体性问题}

关于坐标 $\theta_\magn$ 的整体定义，有一个微妙的问题。我们选用的规范便于讨论 $\theta_\elec$ 的周期性。然而 $A^\sfm$ 中出现的对数表明，这个规范在 $a = 0$ 处是奇异的。此外，沿绕原点一周的延拓 $a \to e^{2 \pi i} a$，$A^\sfm$ 平移 $- \frac{q^2}{2 \pi} d \theta_\elec$。这一平移必须由一个规范变换来补偿：
\begin{equation}
\theta_\magn \to \theta_\magn + q^2 \theta_\elec + C.
\end{equation}
为确定 $C$，我们作一个规范变换，换到新坐标 $\theta_\magn'$：
\begin{equation}
\theta_\magn' = \theta_\magn + \frac{i}{4 \pi} \left( \log \frac{a}{\Lambda} - \log \frac{\bar a}{\bar\Lambda} \right) (q^2 \theta_\elec + C).
\end{equation}
当 $a$ 绕原点一周时，变换后的 $\theta_\magn'$ 是单值的。规范变换后的 $A^\sfm$ 是
\begin{equation}
(A')^\sfm = - \frac{i}{4 \pi} \left( \frac{da}{a} - \frac{d \bar a}{\bar a} \right) (q^2 \theta_\elec + C).
\end{equation}
现在我们关注 $q \theta_\elec = \pi$ 处的行为。这里 $A^\inst = 0$，所以精确的规范场恰由 $(A')^\sfm$ 给出。另一方面，一旦纳入瞬子修正，该点上度量与 $U(1)$ 丛都不再奇异（回忆唯一的奇点出现在 $a=0$、$q \theta_\elec = 2 \pi n$ 处）。又因 $\theta_\magn'$ 是单值的，于是 $(A')^\sfm$ 在这里不可能有奇性——即使延伸到 $a=0$ 也不行（更确切地说，唯一允许的奇性是一条量子化的 Dirac 弦）——这意味着
\begin{equation}
C = - q \pi + 2 \pi  k
\end{equation}
对某整数 $k$ 成立。于是我们的结论是：绕 $a = 0$ 一周时，角坐标的平移为
\begin{subequations}
\begin{align} \label{eq:angle-monodromy}
\theta_\elec &\to \theta_\elec, \\
\theta_\magn &\to \theta_\magn + q^2 \theta_\elec - q \pi. \label{angle-monodromy-varphi}
\end{align}
\end{subequations}
$q^2 \theta_\elec$ 的平移正如环面丛的单值性所预期。$-q \pi$ 的平移则更令人意外，但它正契合我们在第 2.3 节末的讨论：我们那时提出，Wilson 线只有在选定一个局部的二次加细 $\sigma$ 之后才有良好定义。本节至此一直选用「标准的」加细 $\sigma(\gamma_e,\gamma_m) = (-1)^{\gamma_e \gamma_m}$。单值性把 $\gamma_e \to \gamma_e + q^2 \gamma_m$，从而把 $\sigma$ 换成 $\sigma'(\gamma_e,\gamma_m) = (-1)^{q^2 \gamma_m^2} \sigma(\gamma_e,\gamma_m)$。这一加细的改变由 $\theta_\magn$ 平移 $-q \pi$ 来补偿。% TODO 合稿改 \ref{sec:reduction}

\subsection{超凯勒结构}

接下来我们要把 $\CM$ 描述为超凯勒流形。任何 Gibbons–Hawking 度量的超凯勒结构，都由如下三重辛形式决定：
\begin{equation}
\omega^\alpha = d x^\alpha \wedge \left(\frac{d \theta_\magn}{2\pi} + A(\vec{x})\right) + \half \eps^{\alpha \beta \gamma} V dx^\beta \wedge dx^\gamma.
\end{equation}
全纯辛形式 式 (3.3) 则为
\begin{equation}
\hsymp(\zeta) = \frac{1}{4 \pi^2 R} \xi_\magn \wedge \xi_\elec
\end{equation}
其中
\begin{align}
\xi_\magn &= i d\theta_\magn + 2 \pi i A(\vec{x}) + \pi i V \left(\frac{1}{\zeta} da - \zeta d\bar{a}\right), \label{eq:xi-magn} \\
\xi_\elec &= i d\theta_\elec + \pi R \left(\frac{1}{\zeta} da + \zeta d\bar{a}\right). \label{eq:xi-elec}
\end{align}
特别地，由此可知 $\xi_\elec$ 与 $\xi_\magn$ 都是 $(1,0)$ 型的。% TODO 合稿改 \eqref{eq:Omegazeta}

此外，$\xi_\elec$ 可以写成
\begin{equation}
\xi_\elec = \frac{d\CX_\elec}{\CX_\elec}
\end{equation}
其中
\begin{equation} \label{eq:X1-def}
\CX_\elec = \exp \left[ \pi R \frac{a}{\zeta} + i \theta_\elec + \pi R \zeta \bar{a} \right].
\end{equation}
所以 $\CX_\elec$ 是 $\CM$ 上复结构 $J^\pz$ 下的全纯函数。注意：若取 $\gamma$ 为单位电荷，则它与 式 (3.10) 给出的半平坦坐标 $\CX^\sfm_{\gamma}$ 一致，因为此时 $Z_\gamma = a$、$\theta_\gamma = \theta_\elec$。换言之，「电」复坐标不受带电荷粒子的瞬子修正影响：
\begin{equation}
\CX_\elec = \CX^\sfm_\elec.
\end{equation}
% TODO 合稿改 \eqref{eq:X-sf}

为完成 $\CM$ 复几何的描述，还应构造第二个「磁」复坐标 $\CX_\magn$，使得
\begin{equation} \label{eq:second-coord}
\hsymp(\zeta) =- \frac{1}{4 \pi^2 R}\frac{d \CX_\elec}{\CX_\elec} \wedge \frac{d \CX_\magn}{\CX_\magn}.
\end{equation}
这样的 $\CX_\magn$ 必形如
\begin{equation} \label{eq:mag-form}
\CX_\magn = e^{i \theta_\magn + \Phi(a, \bar{a},\theta_\elec, \zeta)}.
\end{equation}

构造 $\CX_\magn$ 最显然的做法，是写出 $\CM$ 上的 Cauchy–Riemann 方程，然后寻找形如 \eqref{eq:mag-form} 的特解。下一节我们改用另一种做法：直接给出 $\CX_\magn$ 的一个特解，其形式对后文特别方便；然后不去验证 Cauchy–Riemann 方程，而是直接验证 \eqref{eq:second-coord}。

\subsection{$\CX_\magn$ 的解}
\label{sec:solving-xmagn}

现在特化到我们的 $\CM$。此时
\begin{equation}
\hsymp(\zeta) = \hsymp^\sfm(\zeta) + \hsymp^\inst(\zeta)
\end{equation}
其中
\begin{align}
\hsymp^\sfm(\zeta) &= - \frac{1}{4 \pi^2 R} \xi_\elec \wedge \left[ i d \theta_\magn + 2 \pi i A^\sfm + \pi i V^\sfm \left(\frac{1}{ \zeta} da -  \zeta d\bar{a}\right) \right],  \\
\hsymp^\inst(\zeta) &= - \frac{1}{4 \pi^2 R} \xi_\elec \wedge \left[ 2 \pi i A^\inst +
\pi i V^\inst \left(\frac{1}{\zeta} da -  \zeta d\bar{a}\right)\right]. \label{eq:instanton-corrections-Omega}
\end{align}
若忽略瞬子修正，所需的磁坐标为
\begin{equation} \CX_\magn^\sfm(\zeta)= \exp\left[-i \frac{R q^2}{2 \zeta} \left(a
\log\frac{a}{\Lambda} -a\right) + i \theta_\magn + i \frac{\zeta R q^2}{2} \left(\bar a
\log\frac{\bar a}{\bar \Lambda} - \bar a\right)\right].
\end{equation}
若取 $\gamma$ 为单位磁荷、$Z_\gamma = \frac{q^2}{2 \pi i} (a \log \frac{a}{\Lambda} - a)$、$\theta_\gamma = \theta_\magn$，它便与半平坦几何中全纯坐标 $\CX^\sfm_\gamma$ 的表达式 式 (3.10) 一致。直接计算可以验证
\begin{equation}
\frac{d \CX^\sfm_\magn}{\CX^\sfm_\magn} = \left[ i d \theta_\magn + 2 \pi i A^\sfm + \pi i V^\sfm \left(\frac{1}{\zeta} da -  \zeta d\bar{a}\right) \right] - \frac{i q^2}{4 \pi}\left( \log \frac{a}{\Lambda} - \log \frac{\bar{a}}{\bar{\Lambda}} \right) \frac{d \CX_\elec}{\CX_\elec},
\end{equation}
从而特别地
\begin{equation} \label{eq:Omega-sf-repeat}
\hsymp^\sfm(\zeta) = - \frac{1}{4 \pi^2 R} \frac{d\CX_\elec}{\CX_\elec} \wedge \frac{d\CX_\magn^\sfm}{\CX_\magn^\sfm},
\end{equation}
如所预期。% TODO 合稿改 \eqref{eq:X-sf}

注意 $\CX_\magn^\sfm$ 绕 $a=0$ 有非平凡的单值性：$\log a$ 与 $\log \bar a$ 的单值性，与 \eqref{angle-monodromy-varphi} 中给出的 $e^{i \theta_\magn}$ 的单值性合起来，给出
\begin{equation}\label{eq:mono2}
\CX_\magn^\sfm \to (-1)^q \CX_\elec^{q^2} \CX_\magn^\sfm.
\end{equation}

接下来纳入瞬子修正。如下文将要证明的，满足 \eqref{eq:second-coord} 的所需 $\CX_\magn$ 可由如下积分公式给出：
\begin{equation}\label{eq:FullX2}
\begin{split}
\CX_\magn = \CX_\magn^\sfm \exp\Biggl[ \frac{i q}{4\pi} & \int_{\ell_+} \frac{d \zeta'}{\zeta'}
\frac{\zeta' + \zeta}{\zeta' - \zeta} \log[1-\CX_\elec(\zeta')^q]  \\
- \frac{i q}{4\pi} & \int_{\ell_-}
 \frac{d \zeta'}{\zeta'} \frac{\zeta' + \zeta}{\zeta' - \zeta} \log[1-\CX_\elec(\zeta')^{-q}] \Biggr],
\end{split}
\end{equation}
其中围道 $\ell_\pm$ 取为位于两个半平面
\begin{equation}
\CU_\pm = \left\{ \zeta: \pm {\rm Re} \frac{a}{\zeta}<0 \right\}.
\end{equation}
内、连接 $0$ 与 $\infty$ 的任意路径。\eqref{eq:FullX2} 中的两个积分贡献分别来自绕 $S^1$ 正绕数与负绕数的瞬子修正。

本节余下的部分验证 \eqref{eq:FullX2} 确实正确。这归结为验证
\begin{equation} \label{eq:tomatch}
- \frac{1}{4 \pi^2 R}\frac{d \CX_\elec}{\CX_\elec} \wedge \frac{d \CX_\magn}{\CX_\magn} = \hsymp(\zeta).
\end{equation}
由 \eqref{eq:FullX2} 有
\begin{equation} \label{eq:dx2x2}
\frac{d \CX_\magn}{\CX_\magn} = \frac{d \CX^\sfm_\magn}{\CX^\sfm_\magn} + \CI_+ + \CI_-
\end{equation}
其中
\begin{equation} \label{eq:ipm}
\CI_\pm = -\frac{i q^2}{4\pi} \int_{\ell_\pm} \frac{d \zeta'}{\zeta'}
\frac{\zeta' + \zeta}{\zeta' - \zeta} \left[ \frac{\CX_\elec(\zeta')^{\pm q}}{1-\CX_\elec(\zeta')^{\pm q}} \frac{d\CX_\elec(\zeta')}{\CX_\elec(\zeta')}\right].
\end{equation}
（这里用到了这样的事实：\eqref{eq:FullX2} 中的积分只通过 $\CX_\elec(\zeta')$ 依赖于 $a,\bar a,\theta_\magn,\theta_\elec$，并且绝对收敛，因此我们可以自由地把微分 $d$ 移入积分号内。）综合 \eqref{eq:tomatch}、\eqref{eq:dx2x2} 与 \eqref{eq:Omega-sf-repeat} 可见，积分 $\CI_\pm$ 必须给出 \eqref{eq:tomatch} 右端的瞬子部分 $\hsymp^\inst(\zeta)$，也就是说我们需要
\begin{equation} \label{eq:intermediate}
\frac{d \CX_\elec(\zeta)}{\CX_\elec(\zeta)} \wedge (\CI_+ + \CI_-) = \frac{d \CX_\elec(\zeta)}{\CX_\elec(\zeta)} \wedge \left[ 2 \pi i A^\inst + \pi i V^\inst \left(\frac{1}{\zeta} da - \zeta d\bar{a}\right)\right].
\end{equation}
为验证这一点，先注意
\begin{equation}
\frac{d \CX_\elec(\zeta)}{\CX_\elec(\zeta)} \wedge \CI_\pm = -\frac{i q^2}{4 \pi} \int_{\ell_\pm} \frac{d \zeta'}{\zeta'}
 \left( \frac{\zeta' + \zeta}{\zeta' - \zeta} \frac{d \CX_\elec(\zeta)}{\CX_\elec(\zeta)} \wedge \frac{d\CX_\elec(\zeta')}{\CX_\elec(\zeta')} \right) \left[ \frac{\CX_\elec(\zeta')^{\pm q}}{1-\CX_\elec(\zeta')^{\pm q}} \right]
\end{equation}
而这里出现的二形式可以改写为
\begin{align}\label{eq:rearrange}
\frac{\zeta' + \zeta}{\zeta' - \zeta} \frac{d \CX_\elec(\zeta)}{\CX_\elec(\zeta)} \wedge \frac{d \CX_\elec(\zeta')}{\CX_\elec(\zeta')} &=
\frac{\zeta' + \zeta}{\zeta' - \zeta} \frac{d \CX_\elec(\zeta)}{\CX_\elec(\zeta)} \wedge \left[\frac{d \CX_\elec(\zeta')}{\CX_\elec(\zeta')}-\frac{d \CX_\elec(\zeta)}{\CX_\elec(\zeta)}\right]\\
&=  - \pi R \frac{d \CX_\elec(\zeta)}{\CX_\elec(\zeta)} \wedge \left[\left(\frac{1}{\zeta'} + \frac{1}{\zeta}\right) da - (\zeta'+\zeta)d\bar a\right],
\end{align}
这里用到了 $\CX_\elec$ 的显式形式 \eqref{eq:X1-def}。于是 \eqref{eq:intermediate} 的左端成为
\begin{equation} \label{eq:intermediate-2}
\begin{split}
\frac{i q^2 R}{4} \frac{d \CX_\elec(\zeta)}{\CX_\elec(\zeta)} \wedge
\biggl( & \int_{\ell_+}  \frac{d \zeta'}{\zeta'}  \left[\left(\frac{1}{\zeta'}+\frac{1}{\zeta}\right)da - ( \zeta'+\zeta)d\bar a\right]\frac{\CX_\elec(\zeta')^{q}}{1-\CX_\elec(\zeta')^{q}}  \\
+ & \int_{\ell_-}  \frac{d \zeta'}{\zeta'}  \left[\left(\frac{1}{\zeta'}+\frac{1}{\zeta}\biggr)da - ( \zeta'+\zeta)d\bar a\right]\frac{\CX_\elec(\zeta')^{-q}}{1-\CX_\elec(\zeta')^{-q}}
\right).
\end{split}
\end{equation}

现在可以计算这些积分了。方便起见，先把两个围道 $\ell_\pm$ 分别形变到恰好位于 $\CU_\pm$ 正中的标准位置，即取
\begin{equation}
\ell_\pm = \left\{ \zeta: \pm \frac{a}{\zeta} \in \IR_- \right\}.
\end{equation}
先考虑乘 $\zeta$ 或 $\frac{1}{\zeta}$ 的那些项。展开几何级数可得：
\begin{align}
\int_{\ell_+}  \frac{d \zeta'}{\zeta'} \frac{\CX_\elec(\zeta')^{q}}{1-\CX_\elec(\zeta')^{q}} &=
\sum_{n>0} \int_{\ell_+} \frac{d \zeta'}{\zeta'}
\exp \left[ \pi R q n\frac{a}{\zeta'} + i q n\theta_\elec + \pi R q n \zeta' \bar{a} \right] \\
&= \sum_{n>0} 2 e^{i q n\theta_\elec}K_0(2 \pi R q|na|).
\end{align}
\eqref{eq:intermediate-2} 中 $\ell_-$ 上的积分给出 $n<0$ 的类似求和。合计得到：\eqref{eq:intermediate-2} 中乘 $\zeta$ 或 $\frac{1}{\zeta}$ 的项等于
\begin{equation} \label{eq:final-part-1}
\frac{d\CX_\elec}{\CX_\elec} \wedge \left( \frac{i q^2 R}{2} \sum_{n \neq 0} e^{i q n\theta_\elec}K_0(2 \pi R q|na|) \right) \left(\frac{1}{\zeta} da - \zeta d\bar{a}\right) = \frac{d\CX_\elec}{\CX_\elec} \wedge i \pi V^\inst\left(\frac{1}{\zeta} da - \zeta d\bar{a}\right).
\end{equation}
对 \eqref{eq:intermediate-2} 中其余的项，类似地有
\begin{align}
\int_{\ell_+}  \frac{d \zeta'}{\zeta'} \zeta' \frac{\CX_\elec(\zeta')^{q}}{1-\CX_\elec(\zeta')^{q}} &=
\sum_{n>0} \int_{\ell_+} \frac{d \zeta'}{\zeta'}
\zeta' \exp \left[ \pi R q n\frac{a}{\zeta'} + i q n\theta_\elec + \pi R q n \zeta' \bar{a} \right] \\
&= - \sum_{n>0} 2 \frac{\abs{a}}{\bar{a}} e^{i q n\theta_\elec}K_1(2 \pi R q|na|)
\end{align}
以及
\begin{align}
\int_{\ell_+}  \frac{d \zeta'}{\zeta'} \frac{1}{\zeta'} \frac{\CX_\elec(\zeta')^{q}}{1-\CX_\elec(\zeta')^{q}} &=
\sum_{n>0} \int_{\ell_+} \frac{d \zeta'}{\zeta'}
\frac{1}{\zeta'} \exp \left[ \pi R q n\frac{a}{\zeta'} + i q n\theta_\elec + \pi R q n \zeta' \bar{a} \right] \\
&= - \sum_{n>0} 2 \frac{\abs{a}}{a} e^{i q n\theta_\elec}K_{1}(2 \pi R q|na|).
\end{align}
把这些与来自 $\ell_-$ 上积分的对应项（它们多带一个负号）合并，我们看到 \eqref{eq:intermediate-2} 中的这些项等于
\begin{equation} \label{eq:final-part-2}
\frac{d\CX_\elec}{\CX_\elec} \wedge \left( - \frac{i q^2 R}{2} \sum_{n \neq 0} e^{i q n\theta_\elec} (\sgn n) \abs{a} K_1(2 \pi R q|na|) \right) \left( \frac{da}{a} - \frac{d \bar a}{\bar a} \right) = \frac{d\CX_\elec}{\CX_\elec} \wedge 2 \pi i A^\inst.
\end{equation}
于是最后，把 \eqref{eq:final-part-1} 与 \eqref{eq:final-part-2} 相加，便如所愿得到 \eqref{eq:intermediate}：对 $\CX_\magn$ 中的围道积分求导，确实正确地产生了瞬子修正 $V^\inst$ 与 $A^\inst$。这也就完成了「$\CX_\magn$ 是 $\CM$ 上所需的『磁』复坐标」这一验证。

\bigskip
\textbf{注}：$\CX_\magn$ 与量子可积系统理论中所谓的「$Q$ 函数」密切相关。\footnote{我们感谢 S. Lukyanov 与我们分享他关于这些函数的笔记。}我们觉得这并非巧合，而是指向与可积场论之间的某种深刻联系。下文的关键方程 式 (5.14) 是热力学贝特拟设的一种形式（见附录 E），这一点更强化了上述感觉。% TODO 合稿改 \eqref{eq:X-integral-mult-explicit}
% TODO 合稿改 \ref{app:TBA}

\subsection{解析性质}

现在考察序对 $(\CX_\magn, \CX_\elec)$ 在 $\zeta$ 平面上的解析行为。

对 $\CX_\elec$ 来说事情很简单：它在 $\zeta \in \IC^\times$ 上解析，在 $\zeta=0,\infty$ 处有本性奇点。对 $\CX_\magn^\sfm$ 也一样，但对完整的 $\CX_\magn$，情况就更复杂：\eqref{eq:FullX2} 中的积分只在离开围道 $\ell_\pm$ 的地方对 $\zeta$ 解析。当 $\zeta$ 跨过这两条围道中的任何一条时，被积函数的极点会越过积分路径。因此，我们关于 $\CX_\magn$ 的表达式定义的是一个分段解析函数，其不连续量由极点的留数决定。引入记号 $(\CX_\magn)_{\ell_+}^{+}$、$(\CX_\magn)_{\ell_+}^{-}$，分别表示 $\zeta$ 沿顺时针或逆时针方向趋近 $\ell_+$ 时 $\CX_\magn$ 的极限，对 $\ell_-$ 也类似。于是不连续量为
\begin{subequations}  \label{eq:X2-disc}
\begin{align}
(\CX_\magn)_{\ell_+}^{+} &= (\CX_\magn)_{\ell_+}^{-} (1 - \CX_\elec^q)^{-q}, \label{eq:X2-disc-a} \\
(\CX_\magn)_{\ell_-}^{+} &= (\CX_\magn)_{\ell_-}^{-} (1 - \CX_\elec^{-q})^q.
\end{align}
\end{subequations}

这些不连续量在下文将起关键作用：我们确实要把它们等同于Kontsevich–Soibelman辛同胚，具体如下。我们把这对复函数 $(\CX_\magn, \CX_\elec)$ 看作给出一个映射
\begin{equation}
\CX: \CM_a \to T_a
\end{equation}
从由 $(\theta_\magn, \theta_\elec)$ 坐标化的实 2-环面 $\CM_a$ 到由 $(X_\magn, X_\elec)$ 坐标化的复化 2-环面 $T_a$。映射 $\CX$ 作为 $\zeta$（以及 $a, \bar{a}, R$）的函数而变化。在第 2.2 节中，我们把 KS 因子 $\CK_\gamma$ 作为 $T_a$ 的辛同胚引入。我们的不连续量 \eqref{eq:X2-disc} 说的是：在射线 $\ell_\pm$ 上，$\CX^{+}$ 与 $\CX^{-}$ 相差一个与 $\CK_{0,\pm q}$ 的复合。% TODO 合稿改 \ref{sec:ReviewKS}

这里出现了一个有趣的现象。考虑 $\CX_\magn$ 在 $a$ 平面中绕 $a=0$ 的单值性。这一单值性有两个来源：\eqref{eq:mono2} 给出的 $\CX_\magn^\sfm$ 的单值性，以及来自 \eqref{eq:X2-disc} 的贡献。这两个贡献实际上互相抵消！这一事实本质上与「半平坦度量在 $a=0$ 处的奇性已被光滑化」有关。另一方面，若把 $\CX_\magn$ 沿 $\zeta = 0$ 解析延拓，它并不回到自身。这一单值性并不带来任何问题。特别地，$\hsymp(\zeta)$ 在 $\zeta = 0$ 附近的行为良好：它只有单极点，正如第 3.1 节的讨论所预期。% TODO 合稿改 \ref{sec:hk-general}

现在考察 $\zeta \to 0,\infty$ 时 $\CX_\elec, \CX_\magn$ 的渐近。$\CX_\elec$ 的渐近可直接由 \eqref{eq:X1-def} 读出：
\begin{equation} \label{eq:X1-asymptotics}
\CX_\elec \sim \begin{cases}
               \exp \left[ \pi R \frac{a}{\zeta} + i \theta_\elec \right] & \text{as } \zeta \to 0, \\
               \exp \left[ \pi R \zeta\bar{a} + i \theta_\elec \right] & \text{as } \zeta \to \infty.
               \end{cases}
\end{equation}
对 $\CX_\magn$，渐近更有意思。当 $\zeta \to 0,\infty$ 时，\eqref{eq:FullX2} 的被积函数简化：有理函数只归结为 $\pm 1$。再展开对数并算出积分，便得
\begin{equation} \label{eq:X2-asymptotics}
\CX_\magn \sim \begin{cases}
               \exp \left[ - i \frac{R q^2}{2\zeta} (a \log (a/\Lambda) - a) + i \theta_\magn + \frac{q}{2\pi i} \sum_{s \neq 0} \frac{1}{s} e^{i s q \theta_\elec} K_0(2 \pi R q \abs{sa}) \right] & \text{as } \zeta \to 0, \\
               \exp \left[ i \frac{\zeta R q^2}{2} (\bar{a} \log (\bar{a}/\bar{\Lambda}) - \bar{a}) + i \theta_\magn - \frac{q}{2\pi i} \sum_{s \neq 0} \frac{1}{s} e^{i s q \theta_\elec} K_0(2 \pi R q \abs{sa}) \right]  & \text{as } \zeta \to \infty.
               \end{cases}
\end{equation}
这些渐近对 $\zeta$ 的一切相位都成立。\eqref{eq:X2-disc} 沿 $\ell_\pm$ 的不连续量不会导致渐近的不连续，因为沿 $\ell_\pm$ 令 $\zeta \to 0,\infty$ 时，跳变指数地接近 $1$：沿 $\ell_+$ 有 $\CX_\elec \to 0$ 指数地快，沿 $\ell_-$ 则 $\CX_\elec^{-1} \to 0$ 指数地快。

另一方面，我们也可以定义另一个函数 $\CX'_\magn$：把 $\CX_\magn$ 顺时针跨过 $\ell_+$ 作解析延拓。由 \eqref{eq:X2-disc-a} 可知，在 $\ell_+$ 的顺时针一侧有
\begin{equation} \label{eq:continued-vs-jumping}
\CX'_\magn = \CX_\magn (1 - \CX_\elec^q)^{q}.
\end{equation}
现在假设把 $\CX'_\magn$ 继续延拓：顺时针延拓到 $\CU_+$ 的边界，再跨入 $\CU_-$。在 $\CU_-$ 中，当 $\zeta \to 0$ 时 $\CX_\elec$ 指数地大。于是由 \eqref{eq:continued-vs-jumping} 可知，$\CX'_\magn$ 与 $\CX_\magn$ 在 $\zeta \to 0$ 时的渐近不同；特别地，$\CX'_\magn$ 并\emph{不}满足 \eqref{eq:X2-asymptotics}。\emph{所以，函数 $\CX_\magn$ 的解析延拓的渐近，并不等于渐近的解析延拓。}这正是斯托克斯现象的标志。

总而言之，我们被引向考虑一个映射 $\CX: \CM_a \to T_a$：它全纯地依赖于 $\zeta$，在 $\zeta \to 0, \infty$ 处呈现斯托克斯现象，且斯托克斯因子由作用在 $T_a$ 上的 KS 辛同胚的复合给出。本文的关键想法是：这一图景对一般的规范理论也成立，而不只是我们这里考虑的阿贝尔理论；事实上，正如我们将在第 5 节看到的，它自动纳入来自相互非局域粒子的多瞬子效应，并给出 $\CM$ 上的精确度量。% TODO 合稿改 \ref{sec:constructing-moduli}

\subsection{微分方程}

上文我们看到，$\CM$ 的超凯勒几何自然地由一个呈现斯托克斯现象的映射 $\CX$ 来描述。斯托克斯现象通常出现在具非正则奇点的线性常微分方程理论中。我们的情形中确实有这样一个微分方程
\begin{equation} \label{eq:de-zeta}
\zeta \partial_\zeta \CX = \CA_\zeta \CX,
\end{equation}
它带有非正则奇点。本节就来识别这个方程。事实上，我们同时还会找到一个伴随方程，它控制对 $S^1$ 半径的依赖：
\begin{equation} \label{eq:de-R}
R \partial_R \CX = \CA_R \CX.
\end{equation}

形如 \eqref{eq:de-zeta}、\eqref{eq:de-R} 的方程，在 $\CA_\zeta$、$\CA_R$、$\CX$ 为有限维矩阵时很常见。此时 $\CA_\zeta$ 与 $\CA_R$ 以从左作矩阵乘的方式作用于 $\CX$，而斯托克斯因子从右作用，所以这两种作用特别地互相对易。我们的情形中，解 $\CX$ 是映射 $\CM_a \to T_a$。斯托克斯因子作为 $T_a$ 的微分同胚作用。$\CA_\zeta$ 与 $\CA_R$ 则作为 $\CM_a$ 的无穷小微分同胚作用，即作为 $(\theta_\magn, \theta_\elec)$ 中的微分算子作用。这两种作用彼此对易，因为它们作用在不同的空间上。

那么，所需方程的来源是什么？它们应当与 $(\CM, g)$ 的某些对称性有关。乍看之下，$(\CM, g)$ 似乎有一个 $U(1)$ 对称性，它只是把 $a \mapsto e^{i \theta} a$。这样的对称性会有显然的物理来源：它来自 $d=4$ 理论的 $U(1)_R$ 对称性。然而我们知道，一旦纳入物质超多重态，这个对称性实际上是反常的。确实，\eqref{eq:A0} 中的 $A^\sfm$ 含有因子 $\log(a/\Lambda)$，而它只在 $a$ 与 $\Lambda$ 的\emph{同时}旋转下才不变。因此这一同时旋转保持度量不变。它并不保持超凯勒形式 $\vec{\omega}$，而是把 $\omega_1$ 与 $\omega_2$ 互相旋转；因此，若与 $\zeta \mapsto e^{i \theta} \zeta$ 作用合起来，它就保持 $\hsymp(\zeta)$ 不变。逐项检查可知，$\CX_\elec$ 与 $\CX_\magn$ 在 $a$、$\Lambda$ 与 $\zeta$ 的这一联合旋转下都不变，这就给出一个微分方程：
\begin{equation} \label{eq:R-symmetry}
\plog{\zeta} \CX = \left( -\plog{\Lambda} + \plog{\bar\Lambda} - \plog{a} + \plog{\bar a} \right) \CX.
\end{equation}
类似地，$d=4$ 理论的反常标度不变性给出一个把 $R$、$a$ 与 $\Lambda$ 同时标度的对称性：
\begin{equation} \label{eq:conformal-symmetry}
\plog{R} \CX = \left( \plog{\Lambda} + \plog{\bar\Lambda} + \plog{a} + \plog{\bar a} \right) \CX.
\end{equation}
这些方程还不是所要的形式 \eqref{eq:de-zeta}、\eqref{eq:de-R}，因为它们仍含有对参数 $(\Lambda, \bar\Lambda, a, \bar a)$ 的导数。所以让我们考察对这些参数的依赖。

$\CX$ 对 $(a, \bar{a})$ 的依赖，完全由它对 $(\theta_\elec, \theta_\magn)$ 的依赖决定，因为要求 $(\CX_\elec, \CX_\magn)$ 在复结构 $J^\pz$ 下全纯。的确，用 \eqref{eq:xi-magn}、\eqref{eq:xi-elec} 作为 $(T^*)^{1,0} \CM$ 的基，我们看到 $\CM$ 上的 Cauchy–Riemann 方程就是
\begin{align}
\partial_a \CX &= \CA_a \CX, \label{eq:de-cr-1} \\
\partial_{\bar a} \CX &= \CA_{\bar a} \CX, \label{eq:de-cr-2}
\end{align}
其中联络形式 $\CA$ 由
\begin{align}
\CA_a &= \frac{1}{\zeta} \left[ - i \pi R \partial_{\theta_\elec} + \pi (V + 2 \pi i R A_{\theta_\elec}) \partial_{\theta_\magn} \right] + 2 \pi A_a \partial_{\theta_\magn}, \\
\CA_{\bar a} &= 2 \pi A_{\bar a} \partial_{\theta_\magn} - \zeta \left[ i \pi R \partial_{\theta_\elec} + \pi (V - 2 \pi i R A_{\theta_\elec}) \partial_{\theta_\magn} \right].
\end{align}
定义。我们也可以类似地处理对 $(\Lambda, \bar\Lambda)$ 的依赖。先注意 $\CX_\elec$ 根本不依赖 $(\Lambda, \bar\Lambda)$。对 $\CX_\magn$ 有 $\Lambda \frac{\partial \CX_\magn}{\partial \Lambda} = \frac{i R q^2 a}{2 \zeta} \CX_\magn$，对 $\bar\Lambda$ 类似。于是，记
\begin{equation}
\CA_\Lambda = \frac{q^2 R a}{2 \zeta} \partial_{\theta_\magn}, \quad \CA_{\bar\Lambda} = \frac{ \zeta q^2 R \bar{a}}{2} \partial_{\theta_\magn},
\end{equation}
就有
\begin{align}
\Lambda \partial_\Lambda \CX &= \CA_\Lambda \CX, \\
\bar\Lambda \partial_{\bar\Lambda} \CX &= \CA_{\bar\Lambda} \CX.
\end{align}
现在可以把方程 \eqref{eq:R-symmetry}、\eqref{eq:conformal-symmetry} 改写为所要的形式 \eqref{eq:de-zeta}、\eqref{eq:de-R}，其中
\begin{align}
\CA_\zeta &= -a \CA_a + \bar{a} \CA_{\bar a} - \Lambda \CA_{\Lambda} + \bar\Lambda \CA_{\bar\Lambda}, \label{eq:zeta-connection-abelian} \\
\CA_R &= a \CA_a + \bar{a} \CA_{\bar a} + \Lambda \CA_{\Lambda} + \bar\Lambda \CA_{\bar\Lambda}.
\end{align}
现在到了关键之处：\eqref{eq:zeta-connection-abelian} 给出的 $\CA_\zeta$ 对 $\zeta$ 的依赖非常简单——它在 $\zeta = 0, \infty$ 处只有单极点：
\begin{equation}
\CA_\zeta = \frac{1}{\zeta} \CA^{(-1)}_\zeta + \CA^{(0)}_\zeta + \zeta \CA^{(1)}_\zeta.
\end{equation}
于是方程 \eqref{eq:de-zeta} 定义了 $\IC\IP^1$ 上的一个亚纯联络，带有两个秩 $1$ 的非正则奇点。这就解释了斯托克斯现象的出现，而我们在上一节已明确看到了它。

一族与 \eqref{eq:de-zeta}、\eqref{eq:de-R}、\eqref{eq:de-cr-1}、\eqref{eq:de-cr-2} 非常相似的微分方程——它们定义了所谓「$tt^*$ 联络」——出现在 \cite{Cecotti:1993rm,MR1213301} 中，背景是 $d=2$ 的 $\CN=(2,2)$ 场论的分析与分类。这种相似并不只是形式上的。特别地，他们对这些方程关于 $\zeta$ 与 $R$ 依赖的解释，同样基于底层场论的 $U(1)_R$ 对称性与标度变换。他们分析中的一个关键点，是联络 $\CA$ 在大 $R$ 时的渐近、斯托克斯因子的显式形式与 $d=2$ 理论中 BPS 态的简并度之间的直接关系。我们的问题中也有类似的关系。事实上，正是这一关系成为理解 KS 穿墙公式的关键。

看一看 \cite{Cecotti:1993rm,MR1213301} 还提示了一个非常有用的技术工具，可用于取得进一步进展：我们应当把这些微分方程转化为 $\CX$ 的黎曼–希尔伯特问题，直接用斯托克斯数据与 $\zeta \to 0, \infty$ 时的渐近来定义。借助这个工具，我们可以立即写下对多个相互非局域 BPS 瞬子的推广。我们将在第 5 节转向这一问题。% TODO 合稿改 \ref{sec:constructing-moduli}

\subsection{更高自旋多重态}

至此我们已相当详细地考察了积分掉单个带电荷超多重态对 $g$ 的修正。同样可以问：单个带电荷的更高自旋多重态——例如含有有质量 $W$ 玻色子的矢量多重态——会带来什么修正？

原则上，这些修正可以通过三维中细致的一圈计算定出。我们改用一个小技巧：考虑 $\CN=4$ 超对称理论的有质量矢量多重态。按 $\CN=2$ 超对称分解，这个多重态含有两个超多重态与一个矢量多重态。另一方面，由于 $\CN=4$ 理论具有更高的超对称，人们预期 $\CM$ 上的度量不会受到任何瞬子修正。理由是：按标准的非重整化定理，Higgs 分支不受修正 \cite{Argyres:1996eh}，而无反常的 $R$ 对称性会把 Higgs 分支与库仑分支混合起来。由此可知，来自 $\CN=2$ 矢量多重态的修正，必定恰好抵消来自两个 $\CN=2$ 超多重态的修正。换言之，至少就这两个 $\CN=2$ 多重态而言，修正是以螺旋度超迹 $\degen(\gamma;u)$ 为权重的。

更一般地，我们可以考虑积分掉任意自旋的 $\CN=2$ 多重态。设 $a_j$ 记第 $j$ 个 $\CN=2$ 多重态（$j=0$ 为超多重态，$j=1$ 为矢量多重态，\dots）对瞬子修正的权重，并归一化为 $a_0 = 1$。上文我们看到 $a_1 = -2$。此外，由「$\CN=4$ 超对称的\emph{任意}多重态的贡献都应当为零」这一事实，得到 $a_{j+2} + 2 a_{j+1} + a_j = 0$。这就定出 $a_j = (-1)^j (j+1)$，所以瞬子修正确实以第二螺旋度超迹为权重。

\subsection{更高秩的推广}
\label{sec:higher-rank}

我们的全部讨论都可以容易地推广到秩 $r$ 的阿贝尔规范理论耦合一组带电荷超多重态的情形。设电荷为 $q_I^{(s)}$，其中 $I=1, \dots, r$ 跑遍各电规范场，$s$ 标记这组超多重态。

Gibbons–Hawking 拟设有 $4r$ 维的推广：底为 $(\IR^3)^r$，纤维为 $(S^1)^r$。底的坐标记作 $(x^{\alpha I}) = \vec x^I$，纤维坐标记作 $\theta_{\magn,I}$，并写
\begin{equation}
g = [V(\vec{x})^{-1}]^{IJ} \left(\frac{d \theta_{\magn,I}}{2\pi} +
A_I(\vec{x})\right)\left(\frac{d \theta_{\magn,J}}{2\pi} +
A_J(\vec{x})\right) + V(\vec{x})_{IJ} d\vec{x}^I d\vec{x}^J,
\end{equation}
其中 $A$ 与 $V$ 由一组微分方程相联系，这组方程说的是
\begin{equation}
\omega^{\alpha} = dx^{\alpha I} \wedge \left( \frac{d\theta_{\magn,I}}{2\pi} + A_I
\right) + \half V_{IJ} \epsilon^{\alpha\beta\gamma} dx^{\beta I } \wedge
dx^{\gamma J}
\end{equation}
对 $\alpha = 1,2,3$ 都是闭 2-形式。

一圈积分给出自然的结果
\begin{equation}
\label{eq:V-sum-r} V_{IJ} =  \mathrm{Im} \tau^0_{IJ} + \sum_{s}
\frac{q^{(s)}_I q^{(s)}_J R}{4 \pi} \sum_{n=-\infty}^\infty \left(
\frac{1}{\sqrt{| q^{(s)}_K R a^K|^2 + (q^{(s)}_K \frac{\theta_\elec^K}{2
\pi} + n)^2}} - \kappa_n \right).
\end{equation}
\eqref{eq:V-sum} 的 Poisson 重求和表明
\begin{equation}
V_{IJ} = V_{IJ}^\sfm + V_{IJ}^\inst,
\end{equation}
其中
\begin{align}
V_{IJ}^\sfm &= \mathrm{Im} \tau^0_{IJ} - \sum_s \frac{q^{(s)}_I
q^{(s)}_J R}{4\pi} \left(\log  \frac{q^{(s)}_K a^K}{\Lambda} + \log \frac{q^{(s)}_K \bar{a}^K}{\bar{\Lambda}} \right), \\
V_{IJ}^\inst &= \sum_s \frac{q^{(s)}_I q^{(s)}_J R}{2\pi}
\sum_{n\neq 0} e^{i n q^{(s)}_K \theta_\elec^K} K_0(2 \pi R \abs{n
q^{(s)}_I a^I}). \label{eq:V-inst-r}
\end{align}
又
\begin{equation}
A_I = A_I^\sfm + A_I^\inst,
\end{equation}
其中
\begin{align}
A_I^\sfm &= \mathrm{Re} \tau^0_{IJ} \frac{d \theta_\elec^J}{2 \pi} + \sum_s
\frac{i q^{(s)}_I
q^{(s)}_J}{8 \pi^2} \left( \log \frac{q^{(s)}_K a^K}{\Lambda} - \log \frac{q^{(s)}_K \bar{a}^K}{\bar{\Lambda}} \right) d\theta_\elec^J, \label{eq:A0-rankr} \\
A_I^\inst &=  - \frac{q^{(s)}_I q^{(s)}_J R}{4 \pi} \left(
\frac{da^J}{q^{(s)}_K a^K} - \frac{d\bar a^J}{q^{(s)}_K \bar a^K}
\right) \sum_{n \neq 0} (\sgn n) e^{i n q^{(s)}_I \theta_\elec^I}
\abs{q^{(s)}_K a^K} K_1(2 \pi R \abs{n q^{(s)}_K a^K}).
\label{eq:A-inst-rankr}
\end{align}
在大 $R$ 时，$V$ 与 $A$ 中的领头项是 $V^\sfm$ 与 $A^\sfm$。只保留这些项，$g$ 就成为半平坦度量，其
\begin{equation}
\tau_{IJ} =\tau^0_{IJ} + \sum_s \frac{q^{(s)}_I q^{(s)}_J}{2 \pi i}
\log \frac{q^{(s)}_K a^K}{\Lambda}.
\end{equation}
这就是 $d=4$ 中积分掉超多重态所得的耦合。

全纯辛形式为
\begin{equation}
\hsymp(\zeta) = - \frac{1}{4 \pi^2 R} \xi^I_e \wedge \xi_{m,I}
\end{equation}
其中
\begin{align}
\xi_e^I &= i d \theta_\elec^I + \pi R \left(\frac{da^I}{\zeta} +
\zeta d \bar a^I \right), \\
\xi_{m,I} &= i d \theta_{\magn,I} + 2 \pi i A_I + i \pi V_{IJ} \left(
\frac{da^J}{\zeta} - \zeta d \bar a^J \right).
\end{align}
如前面一样，电坐标与其半平坦近似一致：
\begin{equation}
\CX^I_e = \exp \left[ \pi R \frac{a^I}{\zeta} + i \theta_\elec^I +
\pi R \zeta \bar a^I \right].
\end{equation}
磁坐标的半平坦近似为
\begin{multline}
\CX^\sfm_{m,I} = \exp \Biggl[ \frac{\pi R}{\zeta} \left(\tau^0_{IJ}
a^J  + \sum_s \frac{q^{(s)}_I}{2 \pi i} q^{(s)}_K a^K \log
\frac{q^{(s)}_K a^K}{e \Lambda}\right) + i \theta_{\magn,I}\,+ \\
\pi R \zeta
\left( \tau^0_{IJ} \bar a^J + \sum_s \frac{q^{(s)}_I}{2 \pi
i} q^{(s)}_K \bar a^K \log \frac{q^{(s)}_K \bar a^K}{e \bar
\Lambda}\right) \Biggr].
\end{multline}
完整的磁坐标由如下积分公式给出：
\begin{equation}\label{eq:FullX2-r}
\begin{split}
\CX_{\magn,I} = \CX_{\magn,I}^\sfm \exp\Biggl[ \sum_s \frac{i
q^{(s)}_I}{4\pi} & \int_{\ell^s_+} \frac{d \zeta'}{\zeta'}
\frac{\zeta' + \zeta}{\zeta' - \zeta} \log[1-\prod_J \CX^J_\elec(\zeta')^{q^{(s)}_J}]  \\
- \frac{i q^{(s)}_I}{4\pi} & \int_{\ell^s_-}
 \frac{d \zeta'}{\zeta'} \frac{\zeta' + \zeta}{\zeta' - \zeta} \log[1-\prod_J \CX^J_\elec(\zeta')^{-q^{(s)}_J}] \Biggr],
\end{split}
\end{equation}
其中 $\ell^s_\pm$ 是位于两个半平面
\begin{equation}
\CU^s_\pm = \left\{ \zeta: \pm {\rm Re} \frac{a^K q^{(s)}_K}{\zeta}<0 \right\}.
\end{equation}
内、连接 $0$ 与 $\infty$ 的任意路径。

注意 $\CX_{m,I}$ 对每个超多重态都有不连续量，如前面一样，它们由 KS 辛同胚 $\CK_{0,q_I^{(s)}}$ 给出。

\begin{yizhu}
本块沿用第 1 章译注已交代的两条约定，此处只作提示：其一，instanton 在本节一律指 $d=4$ 理论中 BPS 粒子的世界线绕 $S^1$ 缠绕所形成的组态，译「BPS 瞬子」，以便与四维瞬子相区别；其二，$\degen(\gamma;u)$ 即原文的 \verb|\degen|，为第二螺旋度超迹，沿用 $\Omega$ 这一符号（不可与 KS 不变量 $\hat\degen(\gamma;u)$ 混同）。
\end{yizhu}
